

\section{SPACY: Spatiotemporal Causal Discovery}

% \textbf{Preliminaries.} A Structural Causal Model \citep{Pearl_2009} (SCM) explicitly defines the causal relationships between variables in the form of functional equations. Formally, an SCM over $D$ variables consists of a 5-tuple $\langle \mathcal{X}, \varepsilon, \mathcal{F}, \mathcal{G}, P(\varepsilon) \rangle$:

% \begin{enumerate}
%     \item Endogenous (observed) variables $\mathcal{X} = \{X^1, X^2, \dots, X^D \}$;
%     \item Exogenous (noise) variables $\varepsilon = \{ \varepsilon^1, \varepsilon^2, \dots, \varepsilon^D \}$ influencing the endogenous variables. 
%     \item A \textit{Directed Acyclic Graph} (DAG) $\mathbf{G}$, denoting the causal links amongst the members of $\mathcal{X}$;
%     \item A set of $D$ functions $\mathcal{F} = \{f^1, f^2, \dots, f^D\}$ determining $\mathcal{X}$ through the equations $X^i = f^i(\text{Pa}_{\mathbf{G}}^i, \varepsilon^i)$, where $\text{Pa}_{\mathbf{G}}^i \subset \mathcal{X}$ denotes the parents of node $i$ in graph $\mathbf{G}$ and $\varepsilon^i \subset \varepsilon$;
%     \item $P(\varepsilon)$, which describes a distribution over noise $\varepsilon$.
% \end{enumerate}

\textbf{{Preliminaries.}} {A Structural Causal Model \citep{Pearl_2009} (SCM) defines the causal relationships between variables in the form of functional equations. Formally, an SCM over $D$ variables \change{in a time-varying system, with time denoted by $t \in [T]$,} consists of a 5-tuple $\displaystyle \langle \mathcal{X}^{(1:T)}, \varepsilon^{(1:T)}, \mathbf{G}, \mathcal{F}, p_{\varepsilon_i^{(1:T)}}(\cdot) \rangle$:}
\change{
\begin{enumerate}[noitemsep, topsep=0pt]
    \item Observed variables {$\mathcal{X}^{(t)} = \{\Xt_1, \dots, \Xt_D \}$};
    \item Noise variables {$\varepsilon^{(t)} = \{ \varepsilon_1^{(t)}, \dots, \varepsilon_D^{(t)} \}$} which influence the observed variables;
    \item A \textit{Directed Acyclic Graph} (DAG) $\mathbf{G}$, which denotes the causal links among the members of $\mathcal{X}^{(1:T)}$ upto a maximum time-lag of $\tau$.
    \item {A set of $D$ functions $\mathcal{F} = \{f_1, \dots, f_D\}$ which determines $\mathcal{X}^{(t)}$ through the equations $\Xt_i = f_i(\text{Pa}_{\mathbf{G}}^i(\leq t), \varepsilon_{i}^{(t)})$, where $\text{Pa}_{\mathbf{G}}^i(\leq t) \subset \mathcal{X}^{(t-\tau:t)}$ denotes the parents of node $i$ in $\mathbf{G}$;}
    \item {$p_{\varepsilon_i^{(t)}}$, which describes a distribution over noise $\varepsilon_i^{(t)}$.}
\end{enumerate}
}
% \begin{figure}[t]
%     \centering
%     \includegraphics[width=\textwidth]{assets/spacy_model.pdf}
%     \caption{Overview of the ELBO calculation for SPACY. The model processes spatiotemporal data $\left\{\mathbf{X}_{1:L}^{(1:T), n}\right\}_{n=1}^N$ to infer latent time series $\left\{\mathbf{Z}_{1:D}^{(1:T), n}\right\}_{n=1}^N$, where $D \ll L$. Causal relationships are modeled using a DAG $\mathbf{G}$ sampled from $q_\phi(\mathbf{G})$. Latent time-series are mapped to grid locations via spatial factors $\mathbf{F}$ sampled from $q_\phi(\mathbf{F})$. Arrows in $\mathbf{G}$ are labeled with edge time-lags.}
%     \label{fig:spacy_model}
% \end{figure}

\begin{figure*}
    \centering
    \includegraphics[width=0.85 \textwidth]{assets/spacy_model.pdf}
    
    \caption{Overview of the ELBO calculation for SPACY. The model processes spatiotemporal data $\left\{\mathbf{X}_{1:L}^{(1:T), n}\right\}_{n=1}^N$ to infer latent time series $\left\{\mathbf{Z}_{1:D}^{(1:T), n}\right\}_{n=1}^N$, where $D \ll L$. Causal relationships are modeled using a DAG $\mathbf{G}$ sampled from $q_\phi(\mathbf{G})$. Latent time-series are mapped to grid locations via spatial factors $\mathbf{F}$ sampled from $q_\phi(\mathbf{F})$. Arrows in $\mathbf{G}$ are labeled with edge time-lags.}
    \label{fig:spacy_model}
\end{figure*}
\textbf{Problem Setting.}
We are given $N$ samples of $L$-dimensional time series with $T$ timesteps each. These $L$ time series are arranged in a $K$-dimensional grid of shape $L_1 \times \ldots \times L_K$ such that $L=\prod_{k=1}^K L_k$, with the spatial coordinates scaled to lie in $\mathcal{G} = [0, 1]^K$. In our setting, we consider $K=2$. We denote the observational time series as $\left\{\mathbf{X}^{(1:T), n}_{1:L} \right\}_{n=1}^{N}$. We assume that the dynamics of the observed data are driven by interactions in a smaller number of \textit{latent} time series. We denote the $D$-dimensional latent time series for each of the $N$ samples as $\left\{ \mathbf{Z}^{(1:T), n}_{1:D} \right\}_{n=1}^N$, with $D << L$. The latent time series is stationary with a maximum time lag of $\tau$, meaning the present is influenced by up to $\tau$ past timesteps. Interactions in the latent time series follow an SCM represented by a DAG $\mathbf{G}$. Our goal is to infer the latent time series $\left\{ \mathbf{Z}^{(1:T), n}_{1:D} \right\}_{n=1}^N$ and the causal graph $\mathbf{G}$ in an unsupervised manner.

\begin{wrapfigure}[12]{r}{0.20\textwidth}
  \begin{center}
  \vspace{-35pt}
    \includegraphics[width=0.20\textwidth]{assets/spacy_pgm.pdf}
  \end{center}
  \caption{Probabilistic graphical model for SPACY. Shaded circles are observed and hollow circles are latent.}
  \label{fig:spacy_pgm_main}
\end{wrapfigure}

% \begin{figure}[b]
%     \centering
%     \begin{overpic}[width=0.60\columnwidth]          {assets/spacy_pgm.pdf}
%     \end{overpic}
    
%     \caption{Probabilistic graphical model for SPACY. Shaded circles are observed and hollow circles are latent.}
%     \label{fig:spacy_pgm_main}
% \end{figure}


\subsection{Forward Model} \label{sec:forward_model} 
We formalize our assumptions about the data generation process using a probabilistic graphical model (Figure \ref{fig:spacy_pgm_main}). We assume that the latent time series $\mathbf{Z}$ is generated by an SCM with causal graph $\mathbf{G}$. The number of latent variables $D$ is a hyperparameter. The spatial correlations between grid points are captured by the spatial factors $\mathbf{F} \in \mathbb{R}^{L \times D}$, which map the latent time series $\mathbf{Z}^{(1:T)}_{1:D} \in \mathbb{R}^{D \times T}$ to the observed time series $\mathbf{X}^{(1:T)}_{1:L} \in \mathbb{R}^{L \times T}$. 
Specifically, the observational time series is generated by applying a grid point-wise non-linearity $g_\ell$, to the product of the spatial factors and latent time series, with additive Gaussian noise. \change{$\ell$ denotes the index of the grid points.} 

\begin{align}
\mathbf{X}_\ell^{(t)} &= g_\ell \left( \left[ \mathbf{F}\mathbf{Z} \right]_\ell^{(t)} \right) + \varepsilon_\ell^{(t)}, \quad
\varepsilon_\ell^{(t)} \sim \mathcal{N}(0, \sigma^2_\ell) \nonumber \\
g_\ell(x) &= \Xi \left( \left[x, \mathscr{E}_\ell \right] \right), \quad \mathscr{E}_\ell \in \mathbb{R}^{f}   \label{eqn:decoder}
\end{align}
We implement the nonlinearity $g_\ell$ as an MLP (multilayer perceptron) $\Xi$ shared across all grid-points, with concatenated embeddings $\mathscr{E} \in \mathbb{R}^{L \times f}$ of dimension $f$.
% The full generative process is described in Equations (\ref{eqn:spacy_1})-(\ref{eqn:spacy_5}).

% \begin{align}
% F_d &= \left[\text{RBF}_d(\nu; \rho_d, \gamma_d)\right] \\
% \rho, \gamma &\sim \mathcal{N}(\mu_{\rho}, \sigma_{\rho}), \mathcal{N}(\mu_{\gamma}, \sigma_{\gamma}, \epsilon^d_t) \\
% Z_d^t &= f_{d}\left(Pa^i_G(<t), Pa^i_G(t)\right) \\
% X &= g(F \times Z + \omega) \\
% \epsilon &\sim \mathcal{N}(0, \sigma^2 I)
% \end{align}

% 
% \begin{wrapfigure}[16]{R}{0.4\textwidth}
%   \begin{center}
%     \includegraphics[width=0.25\textwidth]{assets/spacy_pgm.png}
%   \end{center}
%   \caption{Graphical model representation for Spatial-temporal Causal Discovery (SPACY).
%   Shaded circles are observed and hollow circles are latent.}
%   \label{fig:spacy_pgm}
% \end{wrapfigure}

\begin{figure}[h!]
  \begin{minipage}[b]{0.4\textwidth}
    \centering
    \includegraphics[width=0.8\textwidth]{assets/spacy_pgm.pdf}
    
    
  \end{minipage}
  \hfill
  \begin{minipage}[b]{0.55\textwidth}
    \begin{align*}
    \mathbf{Z}_d^{(t)} &= f_{d}\left(\text{Pa}^d_G(<t), \text{Pa}^d_G(t)\right) + \eta_d^{(t)}\\
    \boldsymbol{\rho}_d &\sim U[0, 1]^K, \boldsymbol{\gamma}_d \sim \mathcal{N}\left(0, I \right)\\
    \mathbf{F}_d &= \left[\text{RBF}_d(x_\ell; \boldsymbol{\rho}_d, \boldsymbol{\gamma}_d)\right]_{\ell=1}^L, \quad x_\ell \in \mathcal{G} \\
    \mathbf{X}_\ell &= g_\ell \left( \left[ \mathbf{FZ} \right]_{\ell} \right) + \varepsilon_\ell \\
    \varepsilon_\ell &\sim \mathcal{N}(0, \sigma^2_\ell I)
    \end{align*}
  \end{minipage}
  \caption{Probabilistic graphical model for SPACY and the generative equations. Shaded circles are observed and hollow circles are latent. }
  \label{fig:spacy_pgm_1}
\end{figure}

% We model the problem as a latent variable inference problem.
% To better understand the proposed hierarchical deep Bayesian model, we divide it into its key components: the spatial factors \( F \), the latent variables \( Z^t \), and the observed variables \( X \). The spatial factors \( F \) capture the geospatial attributes influencing the data, the latent variables \( Z^t \) model the underlying temporal dynamics, and the observed variables \( X \) represent the actual measurements affected by both spatial and temporal factors.


% where \( F \in \mathbb{R}^{L \times D} \) are \( D << L\) spatial factors.  
% We describe each component in more detail.

\textbf{Latent SCM.} We model the latent SCM describing $\mathbf{Z}^{(t)}$ as an additive noise model \citep{hoyer2008nonlinear}:
\begin{equation*}
    \mathbf{Z}_d^{(t)} = f_{d}\left(\text{Pa}^d_\mathbf{G}(<t), \text{Pa}^d_\mathbf{G}(t)\right) + \eta_d^{(t)}, 
\end{equation*}
\change{Note that $\eta_d^{(t)}$ is distinct from the additive Gaussian noise $\varepsilon_\ell^{(t)}$ at the grid-level}. The causal graph $\mathbf{G}$ specifies the causal parents of each node, represented by a temporal adjacency matrix with shape $(\tau+1) \times D \times D$. The parent nodes from previous and current time steps are denoted by $\text{Pa}^d_\mathbf{G}(<t)$ and $\text{Pa}^d_\mathbf{G}(t)$ respectively.  We assume that $\Zt_d$ is influenced by at most $\tau$ preceding time steps, i.e., $\text{Pa}_\mathbf{G}(<t) \subseteq \{ \mathbf{Z}^{(t-1)}, \dots, \mathbf{Z}^{(t-\tau)}\}$. $\mathbf{G}^{(1:\tau)}$ represents the lagged relationships and $\mathbf{G}^{(0)}$ represents the instantaneous edges. The time-lag $\tau$ is treated as a hyperparameter. 

 \change{In principle, SPACY is compatible with any differentiable temporal causal discovery algorithm. In this work, we choose Rhino \citep{gong2022rhino} to model the functional relationships because of its flexibility. It is an identifiable framework that captures both instantaneous effects and history-dependent noise.} We parameterize the structural equations $f_d$ using MLPs $\xi_f$ and $\lambda_f$ shared across all nodes. We use trainable embeddings $\mathcal{E} \in \mathbb{R}^{(\tau+1)\times D \times D}$ with embedding dimension $e$ to distinguish between nodes. $f_d$ is defined as:

\begin{equation}
\resizebox{0.48\textwidth}{!}{$
  f_d\left( \text{Pa}_{\mathbf{G}}(\leq t) \right) = \xi_f \left( \sum_{k=0}^\tau \sum_{j=1}^{D} \mathbf{G}^{(k)}_{j,d} \times \lambda_f \left( \left[ \mathbf{Z}^{(t-k)}_{j}, \mathcal{E}^{k}_{j} \right], \mathcal{E}_{0}^{d} \right) \right)
$},
\label{eqn:rhino_eqn}
\end{equation}
\change{where $\mathbf{G}_{j,d}^{(k)}$ denotes the presence of an edge from node $j$ to $d$ at the $k^\text{th}$ time-lag.}
The noise model is based on conditional spline flows \citep{durkan2019neural}, with the parameters of the spline flow predicted by MLPs $\xi_\eta$ and $\lambda_\eta$, which share a similar architecture to $\xi_f$ and $\lambda_f$. 

\textbf{Spatial Factors.} The low-dimensional latent time series are mapped to the high-dimensional grid by the spatial factors $\mathbf{F} \in \mathbb{R}^{L \times D}$, where the $d^\text{th}$ column represents the influence of the $d^\text{th}$ latent variable on each grid location. To effectively capture the correlation between spatially proximate locations under a single latent variable, we model the spatial factors using kernel functions. In theory, any linearly independent, real analytic family of functions would work (see Section \ref{sec:identifiability}). In practice, we find that radial basis functions (RBFs) work quite well, following \citet{manning2014topographic,sennesh2020neuraltopographicfactoranalysisrbf, 2021farnoosh} (see Appendix \ref{sec:ablation_study} for an ablation study). RBFs not only ensure locality, they are also smooth functions that are parameter-efficient. We assume a uniform prior over the grid $\mathcal{G}$ for the center parameter $\boldsymbol{\rho}_d$ of each kernel, and that the scale parameter $\boldsymbol{\gamma}_d$ comes from a standard normal distribution. Mathematically,
\begin{align}
\boldsymbol{\rho}_d &\sim U[0, 1]^K, \boldsymbol{\gamma}_d \sim \mathcal{N}\left(0, I \right), \label{eqn:f_prior} \nonumber \\
\mathbf{F}_{\ell d} &= \text{RBF}_d(x_\ell; \boldsymbol{\rho}_d, \boldsymbol{\gamma}_d) = \exp\left( -\frac{||x_\ell - \boldsymbol{\rho}_d||^2}{\exp(\boldsymbol{\gamma}_d)} \right),
\end{align}
where $x_\ell$ denotes the spatial coordinates of the $\ell^\text{th}$ grid point.




\subsection{Variational Inference}

Let $\theta$ denote the parameters of the forward model. The likelihood $p_\theta \left( \mathbf{X} \right)$ is intractable due to the presence of latent variables $\mathbf{Z}$, $\mathbf{G}$ and $\mathbf{F}$.  We propose using variational inference, optimizing an evidence lower bound (ELBO) instead.

\begin{proposition}
    The data generation model described in Figure \ref{fig:spacy_pgm_main} admits the following evidence lower bound (ELBO):
    % \begin{align*}
    % & \log p_{\theta}\left( \mathbf{X}^{(1:T), 1:N} \right) \geq
    % \sum_{n=1}^N \Bigg\{ \mathbb{E}_{q_{\phi}(\mathbf{Z}^{(1:T), n}|\mathbf{X}^{(1:T), n}) q_{\phi}(G) q_{\phi}(\mathbf{F})} \Big[ \underbrace{\log p_{\theta}\left( \mathbf{X}^{(1:T), n} | \mathbf{Z}^{(1:T), n}, G, \mathbf{F} \right)}_{\text{Conditional Likelihood}} \\
    % & + \underbrace{\mathbb{E}_{q_{\phi}(\mathbf{Z}^{(1:T), n})} \left[ \log p_{\theta}\left(\mathbf{Z}^{1:T} | G \right) - \log q_{\phi}(\mathbf{Z}^{1:T} | G) \right]}_{\text{Latent Representation}} \Bigg\} + \underbrace{\left[ \log p_{\theta}(G) - \log q_{\phi}(G) \right]}_{\text{Causal Graph}} \\
    % & + \underbrace{\left[ \log p_{\theta}(\mathbf{F}) - \log q_{\phi}(F) \right]}_{\text{Spatial Factors}} \Big]
    % = \text{ELBO}(\theta,\phi)
    % \end{align*}
    \begin{small}
    % \begin{equation}
    \begin{align}
        & \log p_{\theta}\left( \mathbf{X}^{(1:T), 1:N} \right) \geq \sum_{n=1}^N \Bigg\{ \mathbb{E}_{q_{\phi}(\mathbf{Z}^{(1:T), n}|\mathbf{X}^{(1:T), n}) q_{\phi}(\mathbf{G}) q_{\phi}(\mathbf{F})}  \nonumber \\
        & \qquad \Bigg[ {\log p_{\theta}\left( \mathbf{X}^{(1:T), n} | \mathbf{Z}^{(1:T), n}, \mathbf{F} \right)} \nonumber + \Big[ \log p_{\theta}\left(\mathbf{Z}^{(1:T), n} | \mathbf{G} \right) \\
        &\qquad - \log q_{\phi}\left(\mathbf{Z}^{(1:T), n} | \mathbf{X}^{(1:T), n}\right) \Big] \Bigg] \Bigg\} - \text{KL}\left( q_\phi(\mathbf{G}) \mid \mid  p(\mathbf{G})  \right) \nonumber\\
        % & + \mathbb{E}_{q_\phi(\mathbf{G})} {\left[ \log p(\mathbf{G}) - \log q_{\phi}(\mathbf{G}) \right]} \nonumber \\
         & \qquad - \text{KL}\left( q_\phi(\mathbf{F}) \mid \mid  p(\mathbf{F})  \right) = \text{ELBO}(\theta,\phi)     \label{eqn:elbo}
        % & + \mathbb{E}_{q_\phi(\mathbf{F})} {\left[ \log p(\mathbf{F}) - \log q_{\phi}(\mathbf{F}) \right]} 
    \end{align}
    % \end{equation}
    \end{small}

\end{proposition}
See Section \ref{section:elbo_derivation} for the derivation. We outline the computation of the ELBO in Figure \ref{fig:spacy_model}. $q_\phi$ represents the variational distribution parameterized by $\phi$. The first term $\log p_{\theta}(\mathbf{X}^{(1:T), n} | \mathbf{Z}^{(1:T), n}, \mathbf{F})$ in \eqref{eqn:elbo} represents the conditional likelihood of the observed data $\mathbf{X}^{(1:T), n}$ conditioned on $\mathbf{Z}^{(1:T), n}$ and $\mathbf{F}$. The remaining terms represent the KL divergences of the variational distributions from their prior distributions.  Next, we describe the design of the variational distributions, with the full implementation details   in Appendix \ref{app:implementation_details}.

\textbf{Causal graph $q_\phi(\mathbf{G})$.} The variational distribution of the causal graph is modeled as a product of independent Bernoulli distributions, indicating the presence or absence of an edge.
% \vspace{-5mm}
\begin{equation*}
    q_\phi \left( \mathbf{G} \right) = \prod_{k=0}^{\tau} \prod_{i,j=1}^D \text{Bernoulli}\left( \mathcal{W}_{i, j}^{k} \right),
\end{equation*}
where $\mathcal{W} \in \mathbb{R}^{(\tau+1) \times D \times D}$ is a learned parameter.
To estimate the expectation over $q_\phi(\mathbf{G})$, we use Monte Carlo sampling by drawing a single graph sample, employing the Gumbel-Softmax trick \citep{jang2016categorical}. \change{We ensure that the learned graph $\mathbf{G}$ is a DAG using the acyclicity constraint from \citet{zheng2018dags}, which is added to the prior $p(\mathbf{G})$.}

\textbf{Spatial Factor $q_\phi(\mathbf{F})$.} We model the variational distributions of the center ${\boldsymbol{\rho}_d}$ and scale  ${\boldsymbol{\gamma}_d}$ as normal distributions with learnable mean and log-variance parameters $(\boldsymbol{\mu}_{\rho_d}, v_{\rho_d}), (\boldsymbol{\mu}_{\gamma_d}, v_{\gamma_d})$. To sample from $q_\phi(\mathbf{F})$, we first sample ${\boldsymbol{\rho}_d}$ and ${\boldsymbol{\gamma}_d}$ using the reparameterization trick \citep{kingma2014}, and then compute the RBF kernel using these parameters. We ensure that the coordinates of the center lie in $[0,1]^K$ by applying the sigmoid function.
\begin{align*}
    {\boldsymbol{\rho}_d} &\sim \mathcal{N} \left(\boldsymbol{\mu}_{\rho_d}, \exp \left( v_{\boldsymbol{\rho}_d} \right)I \right), {\boldsymbol{\gamma}_d} \sim \mathcal{N}\left(\boldsymbol{\mu}_{\boldsymbol{\gamma}_d}, \exp \left( v_{\gamma_d} \right) I\right)
\end{align*}
\begin{equation*}
    \resizebox{0.48\textwidth}{!}{
    $\mathbf{F}_{\ell d} = \text{RBF}_d(x_\ell; \boldsymbol{\rho}_d, \boldsymbol{\gamma}_d) = \exp\left( -\frac{||x_\ell - \text{sigmoid}(\boldsymbol{\rho}_d)||^2}{\exp(\boldsymbol{\gamma}_d)} \right)$}
\end{equation*}
\textbf{Latent Time Series $q_{\phi}(\mathbf{Z}^{(1:T), n}|\mathbf{X}^{(1:T), n})$.} To obtain latents from the observations, we use a neural network encoder. Specifically, we assume  $q_\phi(\mathbf{Z}^{(1:T), n}|\mathbf{X}^{(1:T), n})$ to be a normal distribution whose mean and log-variance are parameterized by MLPs $\zeta_\mu$ and $\zeta_{\sigma^2}$. 
\begin{align*}
    % \mathbf{Z}^{(t), n} \sim \mathcal{N} \left(\zeta_\mu(\mathbf{X}^{(t), n}), \exp \left( \zeta_{\sigma^2}(\mathbf{X}^{(t), n})\right) \right).
\resizebox{0.48\textwidth}{!}{$
      q_{\phi} \left(\mathbf{Z}^{(1:T), n}|\mathbf{X}^{(1:T), n}\right) = \mathcal{N} \left(\zeta_\mu(\mathbf{X}^{(t), n}), \exp \left( \zeta_{\sigma^2}(\mathbf{X}^{(t), n})\right) \right).
$}
\end{align*}
We sample the latents $\mathbf{Z}$ from the variational distribution using the reparameterization trick.
% \begin{align*}
%     \mathbf{Z}^{(t), n} \sim \mathcal{N} \left(\zeta_\mu(\Sigma_{j=1}^V e_j(\mathbf{X}^{(t), n}, E_j)), \exp \left( \zeta_{\sigma^2}(\Sigma_{j=1}^V e_j(\mathbf{X}^{(t), n}, E_j))\right) \right).
% \end{align*}


% Given data samples $\mathbf{X}^{(1:T), 1:N}$, we first The term $\log p_{\theta}(\mathbf{X}^{(1:T), n} | \mathbf{Z}^{(1:T), n}, \mathbf{F})$ represents the conditional likelihood of the observed data $\mathbf{X}^{(1:T), n}$ conditioned on the latent time-series $\mathbf{Z}^{(1:T), n}$ and the spatial factors $\mathbf{F}$. $\displaystyle \log p_{\theta}\left( \mathbf{Z} | \mathbf{G} \right)$ represents the conditional likelihood of the latent $\mathbf{Z}$ under causal graph $\mathbf{G}$, while $\log q_{\phi}\left( Z | G \right)$.   accounts for the prior distribution and entropy of the latent time-series \( Z \), which is conditioned on the causal graph \( G \). The term \( \log P_{\theta}(G) - \log q_{\phi}(G) \) deals with the prior and entropy of the causal graph \( G \), with sparsity prior applied to ensure a sparse adjacency matrix. Lastly, \( \log P_{\theta}(F) - \log q_{\phi}(F) \) handles the prior and entropy of the spatial factors \( F \), of which we assume uniform priors for location and scales of spatial factors. 

% ELBO 1:
% \begin{align*}
%     \log P_{\theta}(X^{1:T}) \geq & \ \mathbb{E}_{q_{\phi}(Z^{1:T}|X^{1:T}) q_{\phi}(G) q_{\phi}(F)} \Big[ \log P_{\theta}(X^{1:T} | Z^{1:T}, G, F) \tag{Likelihood} \\
%     & + \mathbb{E}_{q_{\phi}(Z^{1:T})} \left[ \log P_{\theta}(Z^{1:T} | G) - \log q_{\phi}(Z^{1:T} | G) \right] \tag{Latent Representation} \\
%     & + \left[ \log P_{\theta}(G) - \log q_{\phi}(G) \right] \tag{Causal Graph} \\
%     & + \left[ \log P_{\theta}(F) - \log q_{\phi}(F) \right] \tag{Spatial Factors} \Big] \\
%     & = \text{ELBO}(\theta,\phi)
% \end{align*}

% ELBO 2:

\subsection{Multivariate Extension}
In many applications, it is important to model interactions between different variates on the same grid. For instance, understanding how global temperature patterns influence precipitation requires analyzing multivariate relationships. We extend SPACY to handle such multivariate time series data, where the observational time series $\mathbf{X} \in \mathbb{R}^{V \times L \times T}$ consists of $V \geq 1$ variates. We modify SPACY to learn variate-specific spatial factors $\mathbf{F}^{(v)}$ by specifying the number of nodes $D_v$ for each variate $v$. We then combine latent representations from all variates and perform causal discovery with them. For more implementation details, refer to Appendix \ref{app:multi-variate}.